\documentclass[level=1,teacher]{kidscode}
\newcommand{\lvA}{\faSeedling\ Just starting}
\newcommand{\lvB}{\faShoePrints\ Getting there}
\newcommand{\lvC}{\faCheckCircle\ Got it}
\newcommand{\lvD}{\faRocket\ Expert}
\begin{document}
\kctitle{Final Project Teacher Guide}{My Interactive Art Show: milestones, rubric, reference solution and showcase}
\tableofcontents

\section{The project in one page}
Each student (or pair) builds a Python turtle program that greets the viewer, asks one or two
questions with \texttt{input()}, and draws a picture or pattern that depends on the answers.
It must use six ingredients: \textbf{a question, a variable, an \texttt{if}, a loop, something random,
and a function made with \texttt{def}}. Students present it at the week-12 showcase and explain one part.

\renewcommand{\arraystretch}{1.35}
\begin{tabularx}{\linewidth}{@{}>{\bfseries}p{2.4cm}>{\bfseries}p{3.4cm}X@{}}
\toprule
When & Milestone & What the students do\\
\midrule
Week 10, last 20 min & M1 Plan & Read the brief together, choose an idea, fill in the plan sheet, open the starter and save it with a project name.\\
Week 11, first half & M2 Core drawing works & Turn the plan into a function that draws the main thing; call it in two or more places.\\
Week 11, second half & M3 Surprise and polish & Add the questions, the \texttt{if}, the loop and \texttt{random}; partner playtest (\enquote{two stars and a wish}); add a title.\\
Week 12 & M4 Showcase & Family gallery and 2-minute demos; exit assessment (20 min); certificates.\\
\bottomrule
\end{tabularx}

\begin{materials}
  \item Print one \textbf{Project Brief} per student (8 pages; the plan sheet is page 4).
  \item On each Mac: the folder \texttt{1-elementry/\allowbreak{}project/\allowbreak{}starter/\allowbreak{}} with \texttt{my\_art\_show.py}
        (it is already inside the copied course folder).
  \item On \textbf{your} Mac only: \texttt{1-elementry/\allowbreak{}project/\allowbreak{}reference/\allowbreak{}night\_sky\_postcard.py}.
        Run it once before week 10 (answer \texttt{40} and \texttt{night}, then try \texttt{5} and \texttt{day}) so you can demo it.
  \item Week 12: feedback cards (last page of this guide), certificates (\texttt{extras/\allowbreak{}certificate}), the exit assessment.
\end{materials}

\begin{teachernote}[Why this project]
The project pulls every Level 1 goal together (L1.G1--G9). It is deliberately an art project with
an optional game twist: the group said they love games and visual art, and art has no
\enquote{losing}, so every student can finish something they are proud of. The rubric values three
things equally: \textbf{creativity}, \textbf{working code} and \textbf{being able to explain it}.
\end{teachernote}

\section{Milestone walkthrough}

\subsection{M1 Plan (end of week 10, about 20 minutes)}
\begin{teacherscript}
\enquote{You are going to be artists who paint with code. Your program will say hello, ask the
viewer a question or two, and draw something that changes because of the answers. Let me show
you mine.} (Run the Night Sky Postcard twice: once with \texttt{40} / \texttt{night}, once with
\texttt{5} / \texttt{day}.) \enquote{Same program, different answers, different picture. Now it is
your turn to plan yours.}
\end{teacherscript}
\begin{enumerate}[itemsep=2pt]
  \item Read page 1 and the six ingredients (page 2) aloud together. Ask: \enquote{Where was the
        loop in my postcard? Where was the \texttt{if}?} (Stars; day or night.)
  \item Students read the idea menu and choose or invent an idea (3 minutes). Pairs agree on one.
  \item They fill in the plan sheet. Insist on the \textbf{sketch} and the \textbf{function name}:
        those two make week 11 much easier.
  \item Check each plan (sign \enquote{Teacher OK}). A good plan has: 1--2 questions that are not
        personal, one function with a clear job, and an idea for each ingredient.
  \item At the Macs: open \texttt{my\_art\_show.py} in Thonny, \menu{File > Save as\dots} with a project
        name (no real names, e.g. \texttt{fish\_party.py}), press Run, see the orange square.
\end{enumerate}
\begin{formativecheck}[M1 is done when]
The plan sheet has a sketch, the questions, a function name with its inputs, and a note for each
ingredient, and the renamed file runs. Quick question: \enquote{What will your function draw, and
what does it need to know?} (Good answer: \enquote{A fish; where to put it and what colour.})
\end{formativecheck}

\subsection{M2 Core drawing works (week 11, first 20--25 minutes of build time)}
The goal is one working function that draws the main thing, called in at least two places.
\begin{enumerate}[itemsep=2pt]
  \item Students change the example \texttt{my\_shape} function into their own (rename it, change the drawing).
  \item First draw it once at \texttt{(0, 0)}. Run. Then call it a second time at another spot,
        e.g. \texttt{fish(-200, 100, PURPLE)}. Run.
  \item Show the coordinate picture on page 3 of the brief when they ask \enquote{where is up?}.
\end{enumerate}
\begin{teacherscript}[Ask while you walk around]
\enquote{Show me the line where you \emph{make} your function. Now show me the line where you
\emph{use} it. What happens if you change this number?}
\end{teacherscript}
\begin{formativecheck}[M2 is done when]
The function draws something recognisable in two different places (or colours) and the program runs
without red errors.
\end{formativecheck}

\subsection{M3 Surprise and polish (week 11, second half)}
\begin{enumerate}[itemsep=2pt]
  \item Add the question(s) in Part 4. Numbers need \texttt{int(\dots)}.
  \item Add the \texttt{if} in Part 5: one answer changes a colour, a shape, or what is drawn.
  \item Put the function call inside a \texttt{for} loop in Part 6, with \texttt{random.randint} for the
        places, sizes or the number of turns.
  \item \textbf{Partner playtest (5 minutes):} partners swap seats, run each other's program, try a
        normal answer and a strange one, and give \enquote{two stars and a wish}.
  \item Polish: a title with \texttt{t.write(\dots)}, \texttt{t.hideturtle()}, a goodbye message.
\end{enumerate}
\begin{formativecheck}[M3 is done when]
All six ingredients are ticked in the comment list at the top of the file, and you have seen each one
work. Ask: \enquote{Run it twice. What is different? Why?} (The random part.)
\end{formativecheck}

\subsection{M4 Showcase (week 12)}
See the showcase plan below. Before guests arrive, each student runs their program once and practises
the 2-minute script with a partner.

\section{Unblocking common problems}
Read the red error message \textbf{with} the student (Thonny shows the line number). Ask them to read
it aloud first: half the time they spot the problem themselves.

{\small\renewcommand{\arraystretch}{1.3}
\begin{longtable}{@{}>{\RaggedRight}p{4.1cm}>{\RaggedRight}p{4.4cm}>{\RaggedRight\arraybackslash}p{\dimexpr\linewidth-8.5cm-4\tabcolsep}@{}}
\toprule
What you see & Likely cause & What to say or do\\
\midrule\endhead
\bottomrule\endfoot
The function draws nothing. & The function is defined but never called. & \enquote{A recipe card does nothing until someone cooks it. Where do you \emph{use} it?} Add e.g. \texttt{fish(0, 0, ORANGE)} below the \texttt{def} block, not indented.\\
\midrule
Everything is drawn on top of itself in the middle. & Every call uses the same \texttt{x, y}, or the function ignores its \texttt{x, y}. & Check the function starts with \texttt{t.penup()}, \texttt{t.goto(x, y)}, \texttt{t.pendown()}.\\
\midrule
Long lines across the picture. & The pen was down while moving to a new spot. & Add \texttt{t.penup()} before \texttt{goto} and \texttt{t.pendown()} after.\\
\midrule
\texttt{TypeError: 'str' object cannot be interpreted as an integer} & \texttt{range(answer)} with text from \texttt{input()}. & \enquote{\texttt{input} always gives words. Turn it into a number with \texttt{int(\dots)}.}\\
\midrule
\texttt{ValueError: invalid literal for int()} & The viewer typed a word like \texttt{lots}. & Fine for Level 1: tell viewers to type digits. Expert students can copy the \texttt{isdigit()} check from the reference.\\
\midrule
The \texttt{if} never happens. & The answer was \texttt{Day} or \texttt{day } but the code checks \texttt{"day"}. & Add \texttt{.strip().lower()} after \texttt{input(\dots)}, or print the answer to see exactly what was typed.\\
\midrule
\texttt{IndentationError} or \texttt{expected an indented block} & Missing 4 spaces under \texttt{def}, \texttt{for} or \texttt{if}; or a missing colon. & Show the \enquote{steps inside the loop are pushed in} rule. In Thonny, select lines and press \keys{Tab} or \keys{Shift}+\keys{Tab}.\\
\midrule
\texttt{NameError: name 'fish' is not defined} & The function is used before it is defined, or spelled differently. & Functions go in Part 3, above where they are used. Check spelling and capitals.\\
\midrule
\texttt{ValueError: empty range for randrange()} & \texttt{random.\allowbreak randint(100, 10)}: the smaller number must come first. & Swap the numbers.\\
\midrule
Fill colour spills or looks odd. & \texttt{begin\_fill()} and \texttt{end\_fill()} are around the wrong lines. & Put \texttt{begin\_fill()} right before the shape's first line and \texttt{end\_fill()} right after the last.\\
\midrule
Drawing is very slow. & Many shapes at normal speed. & \texttt{t.speed(0)} near the top. Keep loops under about 200.\\
\midrule
Part of the picture is missing at the edge. & Coordinates outside $-400\ldots400$ or $-300\ldots300$. & Use the coordinate picture on page 3 of the brief.\\
\midrule
The window will not close or Thonny says the program is still running. & \texttt{turtle.done()} keeps the window open on purpose. & Close the turtle window, or press the red \menu{Stop} in Thonny.\\
\midrule
Student is stuck on the idea, not the code. & Too big a plan. & Shrink it: \enquote{What is the ONE thing your picture needs? Let's draw just that.}\\
\end{longtable}}

\section{Full rubric and scoring}
Use this rubric at M4 (you can also use it informally at M2/M3 to give feedback). Each row scores
1--4 points: Just starting = 1, Getting there = 2, Got it = 3, Expert = 4. The three things we value
most (creativity, working code, explaining) each have their own row, so a student whose code is
simple but who explains it beautifully is rewarded.

{\small\renewcommand{\arraystretch}{1.3}
\begin{longtable}{@{}>{\bfseries\RaggedRight}p{2.3cm}*{4}{>{\RaggedRight\arraybackslash}p{\dimexpr(\linewidth-2.3cm-8\tabcolsep)/4}}@{}}
\toprule
Criterion & \lvA\ (1) & \lvB\ (2) & \lvC\ (3) & \lvD\ (4)\\
\midrule\endhead
\bottomrule\endfoot
Creativity (L1.G2, G5) & Example copied with no changes. & Example with new colours, sizes or numbers. & An own idea: own function, own picture, a real surprise. & Original idea with extra details, a twist, or a game element (e.g. \texttt{onkey}).\\
\midrule
Working code (L1.G6, G9) & Does not run yet. & Runs but stops with an error, or parts are missing. & Runs from start to end without errors for normal answers. & Also handles unexpected answers; tidy code with helpful comments.\\
\midrule
Six ingredients (L1.G1--G5) & 1--2 ingredients. & 3--4 ingredients. & All 6 used. & All 6, used meaningfully: function parameters are really used, the \texttt{if} changes the picture, the loop uses the answer.\\
\midrule
Explaining (L1.G7) & Can say what the program does. & Can point to the loop and the function. & Can explain how one part works, line by line, in own words. & Can explain why they chose it and describe a bug they found and fixed.\\
\midrule
Process and teamwork (L1.G7, G8) & Needed help to plan; worked mostly alone even when paired. & Followed the plan with reminders; took turns sometimes. & Met the milestones; swapped driver/navigator; gave kind playtest feedback. & Helped others, used feedback to improve, gave credit for any ideas borrowed.\\
\end{longtable}}

\textbf{Scoring.} Total = sum of the five rows (5--20). Use it to celebrate and to plan, not to rank:
\begin{itemize}[itemsep=1pt]
  \item \textbf{15--20:} Level 1 goals met with room to spare. Point them to the Level 2 placement quiz with confidence.
  \item \textbf{11--14:} Level 1 goals met. Note one row to grow in Level 2.
  \item \textbf{5--10:} Celebrate what works, then use a buffer session (B2 Free-Build and Remix Studio) to finish before the certificate, or start Level 2 with extra support in weeks 1--2.
\end{itemize}
\textbf{Scratch support path:} use the same rubric. Replace \enquote{Python} with Scratch blocks
(\textbf{ask and wait} + \textbf{answer}, \textbf{set/change} a variable, \textbf{if}, \textbf{repeat},
\textbf{pick random}, a \textbf{My Block}). A Scratch project can score 4 in every row.
\textbf{Pairs:} score each student separately for \enquote{Explaining} and \enquote{Process}; the other rows
are shared.

\section{Reference solution: Night Sky Postcard}
The complete program below is the tested file in \texttt{project/\allowbreak{}reference/\allowbreak{}}. It passes 7 automated
checks (\texttt{test\_night\_sky\_postcard.py}): it asks two questions, draws exactly the number of
stars asked for in random places, uses a black sky at night and a sky-blue sky with a sun by day,
draws four houses with the \texttt{house()} function, and copes with answers like \texttt{lots} or
\texttt{5000}. You do not need to run the checks yourself.

\pyfile{reference/night_sky_postcard.py}

\subsection{Section by section}
Use these notes to explain the example in week 10, or to answer \enquote{how did you do that?} questions.

\begin{description}[leftmargin=0pt,itemsep=5pt,font=\bfseries]
\item[Top comments.] A project starts by saying what it does and which ingredients it uses: good
  habit to copy. The \texttt{SMOKE: run} line is only for the course's automatic window test; students do not need it.
\item[PART 1: Colours.] Named colour variables from the colour-blind-safe Okabe-Ito palette.
  Names like \texttt{ORANGE} are easier to read than \texttt{"\#E69F00"}. Capital letters are a
  Python habit for \enquote{this never changes}.
\item[PART 2: Screen and turtle.] An 800 by 600 window, so $x$ runs from $-400$ to $400$ and $y$
  from $-300$ to $300$. \texttt{speed(0)} draws as fast as possible.
\item[PART 3: Functions.] Two homemade commands. \texttt{box(x, y, width, height, fill\_colour)}
  draws any filled rectangle: the \texttt{for} loop runs twice because a rectangle is \enquote{long
  side, short side} twice. \texttt{house(x, y, colour)} uses \texttt{box()} three times (walls, door,
  windows) and draws a triangle roof with a loop of 3 turns of 120 degrees ($360 \div 3$). This is
  decomposition: a big job made of small named jobs, just like My Blocks in week 5.
\item[PART 4: Greet and ask.] \texttt{print} says hello. The first \texttt{input} asks how many stars.
  The check \texttt{isdigit()} is True only for whole numbers like \texttt{25}, so a word like
  \texttt{lots} does not crash the program: the \texttt{else} picks 30 instead. A second \texttt{if}
  caps the number at 200 so the drawing stays quick. The second \texttt{input} asks day or night;
  \texttt{.strip().lower()} turns \texttt{" Day "} into \texttt{"day"}. Anything else becomes night.
\item[PART 5: Choose colours.] An \texttt{if}/\texttt{else} stores two choices in variables
  (\texttt{sky\_colour}, \texttt{star\_colour}) that later parts use. Day and night also differ by
  \textbf{shape} (sun with rays vs crescent moon), not only by colour.
\item[PART 6: Random stars.] The loop runs \texttt{star\_count} times, the viewer's number. Each time,
  \texttt{random.\allowbreak randint} picks a new place and size, so every postcard is different.
  \texttt{dot(size, colour)} draws a filled circle in one step.
\item[PART 7: Sun or moon.] Another \texttt{if}. The sun is 12 rays (a loop turning 30 degrees,
  $360 \div 12$) plus a filled circle on top. The moon is a yellow circle with a sky-coloured circle
  drawn slightly to the right, leaving a crescent: a neat trick worth showing.
\item[PART 8: Grass and houses.] \texttt{box()} draws the grass. \texttt{house()} is called four times
  with different places and colours: one function, four houses. Ask: \enquote{How would you add a fifth house?}
\item[PART 9: Message.] \texttt{t.write} puts text on the picture. \texttt{turtle.done()} keeps the
  window open until the viewer closes it.
\end{description}

\begin{tip}[Questions to ask about the example]
\enquote{Which line would you change to get 100 stars every time?} (None: you type 100 when it asks.)
\enquote{How do we make the houses purple?} (Change the colour in the four \texttt{house(\dots)} lines.)
\enquote{What happens if we delete \texttt{house(190, -140, WHITE)}?} (Only three houses.)
\enquote{Why do we need \texttt{penup()} before \texttt{goto}?} (Otherwise it draws a line on the way.)
\end{tip}

\section{Showcase plan (week 12)}

\subsection{The family gallery}
Turn the room into a small gallery. Each project gets a \enquote{stand}: a Mac with Thonny open and the
project file loaded, a printed plan sheet beside it, and a card with the project name and the
artist's \textbf{nickname}. Invite the families of every student (the parent notes letter has an
invitation line). 30--40 minutes is plenty.

\begin{center}\small\renewcommand{\arraystretch}{1.3}
\begin{tabularx}{\linewidth}{@{}>{\bfseries}p{2.6cm}X@{}}
\toprule
Role & Job\\
\midrule
Artist & Runs the program for each visitor and gives the 2-minute demo.\\
Assistant & (The partner or the next student.) Types the visitor's answers if the artist wants, resets the program, hands out feedback cards.\\
Visitors & Answer the program's questions, ask one question, fill in a feedback card.\\
Host (you) & Welcomes guests, keeps time (a timer per demo), makes sure every artist gets the same attention, reads the safety reminder.\\
\bottomrule
\end{tabularx}
\end{center}

\subsection{The 2-minute demo script}
Print this or write it on a sticky note for each artist. Practise it once in pairs before guests arrive.
\begin{enumerate}[itemsep=2pt]
  \item \enquote{Hi! My project is called \underline{\hspace{3cm}}. It draws \underline{\hspace{4cm}}.}
  \item \enquote{It will ask you a question. Do you want to answer it?} (Visitor answers; press Run.)
  \item \enquote{Look: because you said \underline{\hspace{2cm}}, it drew \underline{\hspace{3cm}}. That is my \texttt{if}.}
  \item Point to one part of the code: \enquote{This is my function \underline{\hspace{2cm}}. It draws \underline{\hspace{2.5cm}}.
        Here is where I use it inside a loop.}
  \item \enquote{One bug I fixed was \underline{\hspace{4cm}}.}
  \item \enquote{Do you have a question for me?}
\end{enumerate}

\subsection{Feedback cards and certificates}
Visitors use the feedback cards on the last page (\enquote{two stars and a wish}). Read them aloud
to each artist afterwards, then hand out the Level 1 certificate (\texttt{extras/\allowbreak{}certificate}). Write
the student's name on it by hand. A short round of \enquote{one thing I learned, one thing I want to
make next} closes the session.

\begin{safety}[Showcase safety]
\begin{itemize}[itemsep=1pt]
  \item Projects are shown \textbf{on the screen only}. Nothing is uploaded, posted or emailed.
  \item Families may take photos of the artwork \textbf{only} if the artist agrees, and no faces or
        names appear on the screen. Keep photos within the family.
  \item Projects show nicknames only, and never ask visitors for private information.
  \item Credit: if a student borrowed an idea (from a friend, the example, a book), the demo says so.
\end{itemize}
\end{safety}

\begin{teachernote}[Fitting it into the week-12 hour]
A workable plan: 0--10 set up stands and practise the script; 10--35 gallery and demos; 35--55
exit assessment (students only, while families chat or leave); 55--60 certificates and celebration.
If families can only come at the end, do the exit assessment first.
\end{teachernote}

\clearpage
\section{Feedback cards (print and cut)}
\newcommand{\feedbackcard}{%
  \begin{tikzpicture}
    \draw[kcpurple,line width=1.2pt,rounded corners=6pt] (0,0) rectangle (8.2,6.6);
    \node[anchor=north west,font=\bfseries] at (0.3,6.4) {\faPalette\ Feedback for the artist};
    \node[anchor=north west,font=\small] at (0.3,5.75) {Project name:};
    \draw[kcgrey!60] (2.6,5.35) -- (7.9,5.35);
    \node[anchor=north west] at (0.3,4.85) {\faStar\ I liked\dots};
    \draw[kcgrey!60] (0.3,3.95) -- (7.9,3.95);
    \node[anchor=north west] at (0.3,3.65) {\faStar\ I liked\dots};
    \draw[kcgrey!60] (0.3,2.75) -- (7.9,2.75);
    \node[anchor=north west] at (0.3,2.45) {\faMagic\ I wish\dots};
    \draw[kcgrey!60] (0.3,1.55) -- (7.9,1.55);
    \node[anchor=north west,font=\small] at (0.3,1.25) {From (first name or \enquote{a visitor}):};
    \draw[kcgrey!60] (0.3,0.35) -- (7.9,0.35);
  \end{tikzpicture}}
\begin{center}
\feedbackcard\hfill\feedbackcard\par\bigskip
\feedbackcard\hfill\feedbackcard
\end{center}

\end{document}
